\documentclass[blue,fullpaper]{ludusofficial}  % Theme: blue, red, green, purple, orange, cyan
                                              % Article type: fullpaper, letter, poster, survey, shortpaper, artpaper

%% Packages
\usepackage{amsmath}
\usepackage{tabularx}
\usepackage{graphicx} % Ecuaciones muy largas

%% Document metadata
\journalname{}
\journalsubtitle{}
\conferencename{}
\publicationyear{2026}
\institution{}  % This is now set by default in the class file
\journalurl{}  % This is now set by default in the class file
% \articletype{FULL PAPER}  % Set by document class option, or use \articletype{} to customize
%% DOI (optional - leave empty if not assigned yet)
% \articledoi{10.1234/ludus.2026.example}

%% Article information
\title{Diseño e implementación de una CNC Laser controlada por un microcontrolador STM32 Nucleo}
\shorttitle{CNC Laser}  % Short title for page headers
\shortauthor{N.P. \& F.B.R.}  % Short author names for page headers (Separar autores con \&)

\author{%
    \textbf{Nicolas Patricelli}\textsuperscript{1}\\
    [0.2cm]
    \textbf{Federico Barrios Retta}\textsuperscript{2}\\
    [0.2cm]
    \mdseries\small\textsuperscript{1}Facultad de Ingenieria, Universidad Nacional de Cuyo, Argentina\\
    \mdseries\small\textsuperscript{2}Facultad de Ingenieria, Universidad Nacional de Cuyo, Argentina
    % \small Corresponding: john.smith@university.edu
}

\begin{document}


%% Abstract and keywords must be defined BEFORE \maketitle
\begin{abstract}

Abstract

\end{abstract}


%% Maketitle creates full-width header with abstract and keywords, then starts two-column layout
\maketitle

\onecolumn
\tableofcontents
\clearpage
% \twocolumn

\newpage

\section{Introducción y objetivos}

El grabado mediante láser requiere que el desplazamiento
relativo entre la herramienta y la pieza sea coordinado
con el accionamiento del láser. Para ejecutar una trayectoria
definida, deben ser interpretadas las instrucciones de
trabajo y generadas las señales que gobiernan los actuadores
de la máquina.

En el marco de la asignatura \textit{Microcontroladores
y Electrónica de Potencia}, se desarrolló una máquina
de control numérico computarizado (CNC) destinada al
grabado mediante láser. El proyecto permite integrar
los contenidos relacionados con la programación de
sistemas embebidos, la comunicación serie, la utilización
de sensores y el accionamiento de motores paso a paso.

\subsection{Objetivo general}

Diseñar e implementar un sistema de control basado en
un microcontrolador para ejecutar trayectorias de grabado
mediante láser a partir de instrucciones enviadas desde
una computadora.

\subsection{Objetivos específicos}

Se plantean como objetivos específicos la interpretación
de las instrucciones de trabajo, la coordinación de los
desplazamientos en el plano de grabado, el ajuste de la
distancia relativa entre el láser y la pieza, y la regulación
del accionamiento del láser durante la ejecución de las
trayectorias.

Asimismo, se busca establecer una referencia de posición
mediante sensores de fin de carrera y verificar el
funcionamiento del sistema mediante ensayos de movimiento
y grabado.



% ============================================================================= %
\section{Descripción General del Sistema}

La máquina desarrollada consiste en un sistema de control numérico computarizado (CNC) destinado al grabado y al corte mediante láser. Durante estos procesos, el cabezal es desplazado respecto de la pieza siguiendo una trayectoria definida por instrucciones G-code, mientras que el accionamiento del láser es coordinado con el movimiento.

El posicionamiento del cabezal en el plano X–Y se realiza mediante un mecanismo de cinemática CoreXY, accionado por dos motores paso a paso. Se incorpora, además, un eje Z que permite desplazar el cabezal para ajustar la distancia entre el módulo láser y la superficie de trabajo. Este ajuste permite adecuar la posición del cabezal a la altura de la pieza y establecer la distancia de enfoque requerida. El eje Z es accionado de manera independiente de los desplazamientos coordinados en el plano X–Y.

Las trayectorias de trabajo son definidas mediante un archivo G-code generado en una computadora. Las instrucciones son transmitidas al microcontrolador STM32 de una placa NUCLEO-F446RE, donde son interpretadas para generar las señales de paso, dirección y habilitación de los drivers de los motores. Asimismo, se genera una señal de modulación por ancho de pulso (PWM) destinada a comandar el accionamiento del módulo láser.

El sistema incorpora finales de carrera para establecer una referencia de posición mediante un procedimiento de referenciación o \textit{homing}. A partir de dicha referencia, los desplazamientos son determinados mediante las consignas y los pasos comandados a los motores. El posicionamiento se realiza en lazo abierto, sin medición continua de la posición del cabezal.

% [Imagen diagrama de bloques y conexiones] Caption: Armadura Ensamblada de la CNC
% podría representar a la computadora conectada al microcontrolador; desde este último deberían salir las señales hacia los drivers de los motores A, B y Z, y hacia el módulo láser. Los finales de carrera deberían aparecer como entradas al microcontrolador. El bloque de la computadora puede identificarse como «Generación y envío de G-code», dejando pendiente el nombre del programa. Los motores A y B deberían vincularse al mecanismo CoreXY y el motor o conjunto de motores Z al desplazamiento vertical del cabezal. 

% ============================================================================= %
\section{Sistema Mecánico y Actuadores}

\subsection{Armadura de la CNC}

La armadura fue construida siguiendo el diseño
\textit{Core XY}, publicado por Javierarnedo en
Thingiverse \cite{estructura_thingiverse}. Se conservaron
las dimensiones y los perfiles del bastidor original,
sin realizar modificaciones en su configuración
estructural.

El bastidor está constituido por perfiles de aluminio
ranurados de sección nominal de 20~mm $\times$ 20~mm,
dispuestos en dos marcos rectangulares, uno inferior
y otro superior, unidos mediante cuatro montantes
verticales. Esta disposición conforma una estructura
prismática que delimita el espacio de la máquina y
proporciona los puntos de montaje para los componentes
del sistema de movimiento.

Las uniones entre los perfiles se realizan mediante
elementos atornillados y escuadras. Las ranuras
longitudinales permiten fijar los elementos de unión
y los soportes de los componentes a lo largo de los
perfiles, facilitando su posicionamiento durante
el montaje.

El marco superior proporciona soporte al conjunto
de desplazamiento en el plano X--Y, mientras que
los montantes verticales mantienen la separación
entre los marcos y vinculan la parte superior
con la base. De este modo, las cargas transmitidas
por los componentes y las reacciones producidas
durante el movimiento son conducidas hacia
la estructura de apoyo.

Durante el ensamblaje, la alineación de los perfiles
y la perpendicularidad de las uniones resultan
relevantes para conservar la geometría del mecanismo.
Los errores de escuadrado del bastidor pueden
trasladarse al montaje de los elementos de guiado,
mientras que las holguras en las uniones pueden
permitir desplazamientos relativos bajo las
solicitaciones del movimiento.

% Fotografía de los perfiles utilizados.
% Leyenda propuesta: Perfiles de aluminio empleados
% en la construcción del bastidor.

% Fotografía del bastidor ensamblado.
% Leyenda propuesta: Bastidor de la CNC ensamblado.

\subsection{Motores paso a paso}

El accionamiento de los mecanismos de movimiento se
realiza mediante motores paso a paso. Su utilización
permite establecer los desplazamientos a partir de una
secuencia de pasos comandados por el microcontrolador,
mediante los drivers correspondientes.

El posicionamiento se realiza en lazo abierto, sin
realimentación continua de la posición. En consecuencia,
la correspondencia entre los pasos comandados y el
desplazamiento físico depende de que los motores puedan
ejecutar el movimiento sin pérdida de pasos bajo las
condiciones de carga y velocidad previstas.

Se utilizaron cuatro motores paso a paso bipolares
con tamaño de brida NEMA~17. Dos accionan
el mecanismo CoreXY y dos corresponden al eje Z.
Su identificación y sus características eléctricas
se presentan en la Subsección~\ref{subsec:motores-seleccionados}.

% ============================================================================= %
\section{Hardware e Interfaces Eléctricas}

\subsection{Microcontrolador seleccionado}

El sistema de control fue implementado mediante una
placa de desarrollo NUCLEO-F446RE, que incorpora un
microcontrolador STM32F446RET6 de la familia STM32F4.
Este dispositivo dispone de un núcleo Arm Cortex-M4
de 32 bits, con unidad de punto flotante de precisión
simple, 512~KB de memoria Flash y 128~KB de memoria
SRAM.

La unidad de punto flotante proporciona soporte
para las operaciones aritméticas empleadas en el
procesamiento de coordenadas y en los cálculos
de movimiento. Aunque el microcontrolador admite
una frecuencia máxima de 180~MHz, en la configuración
del proyecto se utiliza un reloj de sistema
de 84~MHz.

Para la temporización del movimiento se utilizan
los temporizadores TIM2 y TIM4, asociados a la
generación de pasos de los motores del mecanismo
CoreXY y del eje Z, respectivamente. El canal 1
del temporizador TIM3 se utiliza para generar
la señal PWM que comanda el módulo láser.

La comunicación con la computadora se realiza
mediante USART2, configurada a 115200~bit/s.
La interfaz de puerto serie virtual de la placa
permite establecer esta comunicación a través
de la conexión USB del ST-LINK. Asimismo, este
dispositivo integrado permite programar y depurar
el microcontrolador durante el desarrollo.

Las señales de paso, dirección y habilitación
de los drivers se generan mediante pines GPIO.
También se utilizan entradas digitales para
recibir las señales de los finales de carrera.
Los conectores de expansión de la placa permiten
acceder a los pines necesarios para estas interfaces.

La disponibilidad de temporizadores, comunicación
serie, salidas PWM y pines de entrada y salida
permite concentrar en un mismo microcontrolador
la interpretación de instrucciones, la coordinación
del movimiento y el comando de los actuadores.

\subsection{Motores seleccionados}
\label{subsec:motores-seleccionados}
Para el accionamiento de la máquina se utilizaron
cuatro motores paso a paso bipolares de tamaño
de brida NEMA~17, marca Usongshine, modelo
\texttt{17HS3401}.

Dos motores accionan el mecanismo CoreXY,
correspondientes a los canales A y B. Los otros
dos se utilizan para el desplazamiento del eje Z
los cuales se conectan en paralelo a un mismo driver.

Cada motor dispone de dos bobinas, cuyo accionamiento
requiere invertir el sentido de la corriente
para establecer la secuencia de excitación.
Esta función es realizada por los módulos TB6600,
mientras que el microcontrolador genera las señales
de paso, dirección y habilitación.

La denominación NEMA~17 identifica el tamaño
de la brida de montaje y no determina por sí sola
la corriente nominal, el ángulo de paso ni el par
disponible. Por ello, la caracterización de los
actuadores debe completarse mediante la documentación
correspondiente al modelo utilizado.

\begin{table}[htbp]
    \centering
    \small
    \caption{Identificación y características
    de los motores utilizados.}
    \label{tab:motores-seleccionados}
    \begin{tabular}{ll}
        \hline
        Parámetro & Identificación o valor \\
        \hline
        Marca & Usongshine \\
        Modelo & \texttt{17HS3401} \\
        Tamaño de brida & NEMA~17 \\
        Tipo de accionamiento & Bipolar, dos fases \\
        Cantidad instalada & Cuatro motores \\
        Lote relevado & \texttt{202331120} \\
        Ángulo de paso nominal & 1.8°\\
    \end{tabular}
\end{table}

En el firmware se utiliza un ángulo de paso
de $1{,}8^\circ$, equivalente a 200 pasos completos
por revolución, y una subdivisión de $1/16$ de paso.
Con esta configuración se consideran 3200 pulsos
por revolución. Estos parámetros deben coincidir
con el ángulo nominal del motor y con el ajuste
de micropaso seleccionado en los drivers.

El ajuste de corriente de cada TB6600 debe
establecerse a partir de la corriente nominal
de los bobinados. Para el canal Z debe considerarse,
además, la conexión de dos motores en paralelo,
puesto que la corriente suministrada por el driver
se distribuye entre los bobinados de ambos actuadores.
El valor seleccionado y la configuración de los
interruptores DIP quedan pendientes de documentar.


\subsection{Drivers seleccionados}

Para el accionamiento de los motores paso a paso
se utilizaron módulos driver comercializados como
TB6600. Su función consiste en recibir las señales
de comando generadas por el microcontrolador y
establecer la excitación de los bobinados necesaria
para producir el movimiento.

Mediante estos módulos se separa la generación
de las consignas de movimiento del suministro
de energía a los motores. El microcontrolador
determina la cantidad de pulsos, su frecuencia
y el sentido de desplazamiento, mientras que
el driver gobierna las corrientes de los
bobinados a partir de su alimentación de potencia.

\paragraph{Interfaz de comando.}

El movimiento se comanda mediante las entradas
PUL y DIR. Cada pulso válido recibido en PUL
produce un avance en la secuencia de excitación,
cuya magnitud depende de la subdivisión de pasos
seleccionada. La entrada DIR establece el sentido
de avance de dicha secuencia. Además, se dispone
de una entrada ENA para controlar la habilitación
del accionamiento.

Las entradas de comando se presentan mediante
pares de terminales PUL+/PUL-, DIR+/DIR- y
ENA+/ENA-. Según las especificaciones disponibles,
estas entradas son optoacopladas y están destinadas
a señales de 5~V, con una corriente de señal
comprendida entre 8 y 15~mA. Por lo tanto, su
conexionado debe considerar tanto la tensión
aplicada como la corriente requerida por
la interfaz de entrada.

\paragraph{Configuración de corriente y micropasos.}

El módulo permite seleccionar la corriente
de salida mediante interruptores DIP. Este ajuste
establece el nivel de corriente utilizado para
la excitación de los bobinados y debe determinarse
a partir de las características del motor conectado.
La corriente máxima admitida por el driver no
constituye, por sí sola, el valor que debe utilizarse
en el proyecto.

También puede seleccionarse la subdivisión
de pasos. En el funcionamiento por micropasos,
las corrientes de los bobinados son modificadas
para generar posiciones intermedias entre los
pasos completos del motor. Esto permite incrementar
la cantidad de pulsos requeridos por vuelta
y reducir el incremento angular nominal asociado
a cada pulso.

Para un motor con ángulo de paso completo
$\alpha$, la cantidad de pulsos por vuelta
se expresa como

\begin{equation}
    N_{\mathrm{pulsos/vuelta}}
    = \frac{360^\circ}{\alpha}\,m,
\end{equation}

donde $m$ es la cantidad de micropasos por paso
completo. Para $\alpha=1{,}8^\circ$ y $m=16$,
se requieren 3200 pulsos por vuelta.

El firmware del proyecto considera una subdivisión
de 16 micropasos por paso completo. En consecuencia,
la configuración efectiva de los interruptores
del driver debe coincidir con este valor para
mantener la correspondencia entre los pulsos
comandados y el desplazamiento calculado.

\paragraph{Características eléctricas.}

En la Tabla~\ref{tab:tb6600} se resumen las
especificaciones recopiladas para el módulo.
Estos valores deben contrastarse con la documentación
de la versión instalada. La tensión de alimentación,
la corriente seleccionada y la posición de los
interruptores DIP utilizados en la CNC se
documentarán junto con el esquema de conexiones.

\begin{table}[htbp]
    \centering
    \begin{tabular}{p{0.58\linewidth}p{0.30\linewidth}}
        \hline
        Parámetro & Valor declarado \\
        \hline
        Tensión de alimentación de potencia
            & 9--42~V CC \\
        Intervalo de corriente de salida
            & 0,7--4,0~A \\
        Niveles de corriente seleccionables
            & 8 \\
        Frecuencia máxima de pulsos de entrada
            & Hasta 20~kHz \\
        Tensión de las señales de entrada
            & 5~V \\
        Corriente de las señales de entrada
            & 8--15~mA \\
        Pulsos por vuelta para un motor de $1{,}8^\circ$
            & 200--6400 \\
        Aislamiento de las entradas
            & Optoacoplado \\
        Protecciones declaradas
            & Sobrecorriente y sobretemperatura \\
        \hline
    \end{tabular}
    \caption{Especificaciones recopiladas del módulo
    TB6600, pendientes de verificación documental.
    Debe confirmarse si los valores de corriente
    indicados corresponden a corriente pico o eficaz.}
    \label{tab:tb6600}
\end{table}

La frecuencia de los pulsos de entrada determina
la velocidad angular comandada. Para una configuración
que requiere $N_{\mathrm{pulsos/vuelta}}$ pulsos
por vuelta, la velocidad expresada en revoluciones
por minuto resulta

\begin{equation}
    n = \frac{60 f_{\mathrm{PUL}}}
    {N_{\mathrm{pulsos/vuelta}}}.
\end{equation}

Esta relación establece la conversión entre
frecuencia de comando y velocidad nominal.
El funcionamiento sin pérdida de pasos depende
también del par disponible, la carga, la aceleración
y las condiciones de alimentación del conjunto.

\subsection{Conexiones eléctricas}

Las interconexiones entre la placa de desarrollo,
los drivers de los motores, los finales de carrera
y el módulo de accionamiento del láser se representan
en la Figura~\ref{fig:esquematico-electrico}.
El esquema fue elaborado en KiCad e identifica
los terminales de los módulos, las redes de alimentación
y las señales de control utilizadas por el firmware.

\paragraph{Distribución de la alimentación.}
Se dispone de una fuente con entrada de
$220\,\mathrm{V}$ de corriente alterna y salidas de
$24\,\mathrm{V}$, $12\,\mathrm{V}$ y $5\,\mathrm{V}$ de corriente continua.
La alimentación de potencia de los drivers TB6600
se realiza mediante la salida de $12\,\mathrm{V}$,
con los tres módulos conectados en paralelo.
La placa NUCLEO-F446RE se alimenta mediante
[completar: USB o alimentación externa].
La utilización de la salida de $5\,\mathrm{V}$ y
la alimentación del módulo láser deberán identificarse
de acuerdo con el montaje realizado.

En el esquema se distinguen las alimentaciones
de potencia, las tensiones utilizadas por las
interfaces de control y sus correspondientes retornos.
Las uniones entre referencias eléctricas se representan
según el cableado del sistema.

\paragraph{Interfaz con los drivers de los motores.}
El accionamiento de los motores es comandado mediante
señales de paso, dirección y habilitación.
Las señales de paso y dirección son independientes
para los canales A, B y Z, mientras que la habilitación
es compartida por los tres drivers.

El cableado documentado en el firmware corresponde
a una conexión de ánodo común: los terminales
\texttt{PUL+}, \texttt{DIR+} y \texttt{ENA+}
se conectan a $3{,}3\,\mathrm{V}$, y los terminales
negativos reciben las señales de los GPIO.
La correspondencia entre señales y pines se presenta
en la Tabla~\ref{tab:senales-electricas}.

\paragraph{Conexión de los motores.}
Cada motor bipolar dispone de dos bobinas.
Los terminales \texttt{A+} y \texttt{A-} del driver
se conectan a los extremos de una bobina,
y los terminales \texttt{B+} y \texttt{B-},
a los extremos de la otra.

Los motores A y B (denotados por la numeración 1 y 2 respectivamente, para diferenciarlos de la notación utilizada para representar el par de bobinas de los motores paso a paso) del mecanismo CoreXY disponen
de drivers independientes. Para el eje Z,
el firmware documenta dos motores conectados
en paralelo a un mismo driver. Esta disposición
debe reflejarse mediante la conexión de cada salida
del driver a los terminales equivalentes de ambos
motores.

\paragraph{Interfaz con los finales de carrera.}
Se utilizan tres señales de límite, correspondientes
a X, Y y Z, conectadas a PB0, PB1 y PB2,
respectivamente. Los módulos de final de carrera
se representan mediante sus terminales de alimentación,
masa y señal, conservando la identificación
del conector utilizado.

El firmware configura estas entradas con resistencias
internas de polarización a masa e interpreta
un nivel bajo como la activación del límite.
Por lo tanto, debe comprobarse que la salida
de los módulos presente el nivel de reposo
y de activación esperado por el programa.
La tensión de alimentación y la disposición física
de sus terminales quedan pendientes de confirmación.

\newpage

\paragraph{Interfaz con el accionamiento del láser.}
La señal de control del láser se genera mediante
el canal 1 de TIM3 y se obtiene en PC6,
disponible en el terminal 4 del conector CN10
de la NUCLEO-F446RE. Esta red se identifica
como \texttt{LASER\_PWM}.

Con la configuración del temporizador utilizada
en el firmware, la frecuencia nominal de la señal
es de $1\,\mathrm{kHz}$. Su ciclo de trabajo
es modificado para establecer la consigna
de accionamiento del módulo láser.
La identificación del terminal receptor,
la tensión de alimentación y los niveles eléctricos
admitidos por dicho módulo quedan pendientes
de completar.


\begin{figure}
    \centering
    \includegraphics[width=0.75\linewidth]{attachments/Esquematico.pdf}
    \caption{Esquema de interconexión de la CNC:
    alimentación, control de motores, finales
    de carrera y accionamiento del láser.}
    \label{fig:esquematico-electrico}
\end{figure}

\begin{table}[htbp]
    \centering
    \small
    \caption{Asignación de señales de control y entradas
    de límite en la NUCLEO-F446RE.}
    \label{tab:senales-electricas}
    \begin{tabular}{lll}
        \hline
        Señal & GPIO & Destino o procedencia \\
        \hline
        \texttt{STEP\_A}    & PC0 & \texttt{PUL-} del driver A \\
        \texttt{DIR\_A}     & PC1 & \texttt{DIR-} del driver A \\
        \texttt{STEP\_B}    & PC2 & \texttt{PUL-} del driver B \\
        \texttt{DIR\_B}     & PC3 & \texttt{DIR-} del driver B \\
        \texttt{STEP\_Z}    & PC4 & \texttt{PUL-} del driver Z \\
        \texttt{DIR\_Z}     & PC5 & \texttt{DIR-} del driver Z \\
        \texttt{MOTOR\_EN}  & PC7 & \texttt{ENA-} de los tres drivers \\
        \texttt{LIMIT\_X}   & PB0 & Señal del final de carrera X \\
        \texttt{LIMIT\_Y}   & PB1 & Señal del final de carrera Y \\
        \texttt{LIMIT\_Z}   & PB2 & Señal del final de carrera Z \\
        \texttt{LASER\_PWM} & PC6 & Entrada de control del módulo láser \\
        \hline
    \end{tabular}
\end{table}

% Incorporar aquí el PDF exportado desde KiCad.
% Sustituir la ruta por la del archivo utilizado.

\clearpage

% ============================================================================= %
\section{Arquitectura del Firmware}

\subsection{Organización modular}

El firmware fue desarrollado en lenguaje C mediante
una organización modular. Las responsabilidades
del sistema se distribuyen entre módulos de
interpretación de instrucciones, generación
de movimiento, referenciación, control del láser
y habilitación de los drivers.

Los módulos funcionales se organizan mediante archivos
de encabezado, en los que se declara su interfaz pública,
y archivos de implementación, en los que se definen las
operaciones correspondientes. Los encabezados de
configuración reúnen constantes y asignaciones de pines,
sin requerir necesariamente un archivo de implementación.

El estado interno y las funciones auxiliares de los
módulos se mantienen, cuando corresponde, mediante
declaraciones con el especificador \texttt{static}.
En las declaraciones a nivel de archivo, este especificador
establece enlace interno y evita que sus identificadores
sean utilizados directamente desde otras unidades de
compilación. De esta manera, el acceso al módulo se realiza
mediante las funciones declaradas en su encabezado.

Las variables locales declaradas con \texttt{static}
conservan su valor entre llamadas sucesivas. Este recurso
es utilizado, por ejemplo, para recordar las condiciones
de continuidad entre segmentos de movimiento. El
ocultamiento del estado interno permite modificar su
representación sin alterar la interfaz utilizada por
los demás módulos.

\begin{figure}
    \centering
    \includegraphics[width=0.5\linewidth]{attachments/UML_Diagram.png}
    \caption{Diagrama de clase en formato UML}
\end{figure}

A continuación, se muestra el árbol de la estructura del proyecto:

\begin{figure}[htbp]
    \centering
    \begin{minipage}{0.95\linewidth}
    \small
\begin{verbatim}
Code/
|-- CNC_laser.ioc
|-- Core/
|   |-- Inc/
|   |   |-- main.h
|   |   |-- corexy_config.h
|   |   |-- gcode_parser.h
|   |   |-- stepper.h
|   |   |-- z_axis.h
|   |   |-- homing.h
|   |   |-- laser.h
|   |   |-- motor_enable.h
|   |   `-- ... encabezados de soporte
|   |-- Src/
|   |   |-- main.c
|   |   |-- gcode_parser.c
|   |   |-- stepper.c
|   |   |-- z_axis.c
|   |   |-- homing.c
|   |   |-- laser.c
|   |   |-- motor_enable.c
|   |   |-- stm32f4xx_it.c
|   |   |-- stm32f4xx_hal_msp.c
|   |   `-- ... archivos de soporte
|   `-- Startup/
|-- Drivers/
|   |-- CMSIS/
|   `-- STM32F4xx_HAL_Driver/
|-- STM32F446RETX_FLASH.ld
`-- STM32F446RETX_RAM.ld
\end{verbatim}
    \end{minipage}
    \caption{Estructura seleccionada del proyecto de firmware.
    Se abrevian los archivos de soporte y las bibliotecas.}
    \label{fig:estructura-firmware}
\end{figure}

\paragraph{Coordinación de la aplicación: \texttt{main.c}.}
En este archivo se inicializan los periféricos y los
módulos funcionales, y se implementa el bucle principal.
También se construyen las líneas recibidas, se atienden
los comandos especiales y se determinan las acciones
correspondientes a las instrucciones G-code. Se conservan
el modo de coordenadas, el avance modal, el segmento
X--Y pendiente y las banderas asociadas a la comunicación
y a la alarma.

\paragraph{Interpretación de instrucciones:
\texttt{gcode\_parser}.}
Este módulo transforma una línea de texto en una
estructura \texttt{GcodeCommand}. En ella se identifican
las coordenadas, el avance, el parámetro de accionamiento
del láser y los comandos reconocidos. El módulo realiza
la interpretación de la instrucción; la decisión sobre
su ejecución y la conservación del estado modal
corresponden a la aplicación principal.

\paragraph{Movimiento CoreXY: \texttt{stepper}.}
Se implementa la transformación de las consignas X--Y
en desplazamientos de los motores A y B, junto con la
coordinación de los pasos y el perfil de velocidad.
El estado del movimiento y las posiciones lógicas
se conservan internamente. La generación de pulsos
es atendida mediante interrupciones de TIM2, mientras
que la interfaz permite solicitar desplazamientos,
consultar o establecer la posición lógica y detener
el movimiento.

\paragraph{Movimiento del eje Z: \texttt{z\_axis}.}
Se implementa el accionamiento del eje Z mediante
consignas de posición y avance. El módulo conserva
la posición lógica y el estado del desplazamiento,
y actualiza los pasos mediante interrupciones de TIM4.
Su ejecución es independiente del mecanismo de
almacenamiento y continuidad de segmentos X--Y.

\paragraph{Referenciación: \texttt{homing}.}
Se implementa la secuencia de búsqueda de los finales
de carrera en el orden Z, X e Y. La maniobra comprende
una aproximación inicial, un retroceso, una segunda
aproximación a menor velocidad y una separación final
del contacto. Se utilizan las funciones de movimiento
de los módulos CoreXY y Z, y se mantiene una bandera
que identifica la ejecución de la referenciación.

\paragraph{Accionamiento del láser: \texttt{laser}.}
Se administra la habilitación del láser y la actualización
del ciclo de trabajo de la señal PWM generada por TIM3.
El parámetro S se limita al intervalo configurado y se
convierte en un valor de comparación del temporizador.
Se distingue entre la deshabilitación del láser y la
anulación momentánea de su salida durante un movimiento
rápido.

\paragraph{Habilitación de los drivers:
\texttt{motor\_enable}.}
Se administra la señal común de habilitación de los
tres drivers y se conserva su estado lógico. La interfaz
permite habilitar, deshabilitar y consultar dicho estado,
y es utilizada por la aplicación para atender los
comandos M17 y M84.

\paragraph{Configuración del movimiento:
\texttt{corexy\_config.h}.}
Este encabezado reúne las asignaciones de pines y
los parámetros de conversión entre desplazamiento
y pasos de los motores CoreXY, junto con constantes
del perfil de velocidad y límites del área de trabajo.
Los parámetros correspondientes al eje Z se encuentran
definidos en \texttt{z\_axis.h}.

\subsection{Inicialización y programa principal}

Al comenzar la ejecución se inicializa
la biblioteca HAL, se configura el reloj
del sistema y se establecen los parámetros
de los GPIO, USART2 y los temporizadores
TIM2, TIM3 y TIM4.

Posteriormente se inicializan los módulos
de referenciación, habilitación de motores,
movimiento del eje Z, movimiento CoreXY
y control del láser. También se inicia
la recepción de datos por interrupción
y se transmite el mensaje de bienvenida.

Una vez completada la inicialización,
el programa permanece en un lazo principal.
En él se atienden los mensajes pendientes,
se verifica la condición de ejecución
del movimiento X--Y almacenado y se extraen
los caracteres del buffer de recepción
para construir líneas completas.

Cada línea recibida es procesada mediante
la función \texttt{process\_gcode\_line()}.
Los comandos especiales son atendidos
directamente, mientras que las instrucciones
G-code son analizadas por el módulo
\texttt{gcode\_parser}. La estructura obtenida
es utilizada para determinar las acciones
sobre el movimiento, el láser y la habilitación
de los motores.

\subsection{Coordinación de los módulos}

Las consignas X--Y son calculadas considerando
el modo de coordenadas absoluto o relativo.
El movimiento se almacena temporalmente
en una estructura que contiene el destino,
el avance, la dirección del segmento
y la información asociada al láser.

Cuando se recibe un movimiento posterior,
se compara su dirección con la del segmento
pendiente para decidir las condiciones
de continuidad entre ambos. Si no se recibe
una continuación, el movimiento pendiente
es ejecutado al superar el tiempo de espera
establecido.

La función \texttt{flush\_pending\_move()}
aplica la condición de accionamiento del láser
correspondiente al segmento e invoca
\texttt{stepper\_move\_to()} para ejecutar
el desplazamiento CoreXY.

Las consignas Z son atendidas mediante
\texttt{z\_axis\_move\_to()}, de manera
independiente del mecanismo de almacenamiento
de segmentos X--Y. La referenciación
es solicitada mediante el comando
\texttt{\$H}, que invoca
\texttt{homing\_run()}.

\subsection{Ejecución mediante interrupciones}

La recepción UART se realiza por interrupción,
de a un carácter. Los caracteres ordinarios
son almacenados en un buffer circular para
su procesamiento posterior por el programa
principal. Determinados caracteres de control
activan banderas que solicitan el envío
de mensajes.

Las interrupciones de TIM2 y TIM4 invocan
las rutinas de generación de pasos de los
módulos CoreXY y Z, respectivamente.
En estas rutinas se actualiza el avance
del movimiento y se establece la temporización
del siguiente evento de paso.

Las funciones de movimiento son bloqueantes
para el programa principal: su ejecución
no retorna hasta que finaliza o se interrumpe
el desplazamiento. Durante esa espera,
las interrupciones habilitadas continúan
siendo atendidas. Por lo tanto, la recepción
de caracteres puede mantenerse mientras
se ejecuta un movimiento, aunque el
procesamiento de nuevas líneas queda
condicionado por la disponibilidad
del programa principal y del buffer.

\subsection{Representación de los estados de operación}

El estado de la aplicación se conserva
mediante variables y banderas distribuidas
entre los módulos. Se representan, entre
otras condiciones, la existencia de una
alarma, la ejecución de la referenciación,
la presencia de un movimiento pendiente
y la actividad de los movimientos
CoreXY y Z.

La habilitación de los motores, la habilitación
del láser y el modo de coordenadas constituyen
condiciones adicionales que pueden coexistir
con distintos modos de operación. Por ello,
no deben interpretarse como estados globales
mutuamente excluyentes.

A partir de estas condiciones puede elaborarse
un diagrama funcional de operación que muestre
la inicialización, el procesamiento de comandos,
la ejecución de movimientos, la referenciación
y el bloqueo por alarma. Esta representación
describe el comportamiento del firmware,
cuya implementación se basa en llamadas
a funciones y evaluación de banderas.

% ============================================================================= %
\section{Comunicación e Interpretación de G-Code}

\subsection{Generación y envío de instrucciones mediante LaserGRBL}

La preparación de los trabajos y el envío de instrucciones
a la máquina se realizan mediante LaserGRBL. Este programa
permite importar imágenes, establecer parámetros del trabajo
y generar una secuencia de instrucciones G-code. También
proporciona una interfaz para transmitirlas al controlador
y recibir sus respuestas.

La utilización de una herramienta existente permite
delimitar el desarrollo al sistema embebido encargado
de interpretar las consignas y accionar la máquina.
La implementación de los algoritmos de procesamiento
de imágenes y generación de trayectorias no forma parte
del alcance del proyecto. Sin embargo, las instrucciones
generadas deben pertenecer al conjunto reconocido por
el firmware y utilizar parámetros compatibles con
la configuración de la CNC.

En los trabajos de grabado de imágenes pueden generarse
secuencias de segmentos lineales con diferentes valores
de accionamiento del láser. La ejecución de estas
trayectorias requiere conservar el orden de las instrucciones
y asociar cada segmento con su consigna de velocidad
y de accionamiento. La coordinación mecánica del barrido
se desarrolla en la sección de generación de movimiento.

\subsection{Comunicación serie y tratamiento de mensajes}

La comunicación con la computadora se establece mediante
USART2, configurada a 115200~bit/s, con ocho bits de datos,
un bit de parada y sin paridad. La conexión USB del ST-LINK
permite acceder a esta interfaz mediante un puerto
serie virtual.

La recepción se realiza por interrupción, de a un carácter.
Los caracteres ordinarios son almacenados en un buffer
circular y posteriormente extraídos por el programa
principal. Al encontrarse un retorno de carro o un salto
de línea, se termina la cadena y se procesa la instrucción
construida. Las líneas vacías no son enviadas al parser.

El buffer circular dispone de 256 posiciones, de las
cuales 255 pueden contener caracteres pendientes.
El buffer de construcción de líneas dispone de 96
posiciones, incluyendo el terminador de cadena.
Estos límites condicionan la cantidad de información
que puede conservarse mientras el programa principal
se encuentra ocupado.

Para mantener el intercambio con LaserGRBL se utilizan
respuestas y mensajes con formato compatible con su
interfaz: \texttt{ok}, \texttt{error:n},
\texttt{ALARM:n} y \texttt{[MSG:...]}.
Las instrucciones que alcanzan el final del procesamiento
reciben una respuesta \texttt{ok}; determinados errores
producen una respuesta \texttt{error:n} y un retorno
anticipado.

La respuesta \texttt{ok} no debe interpretarse como
una confirmación de que todos los desplazamientos
asociados ya finalizaron. En particular, una consigna
X--Y puede haber quedado almacenada como movimiento
pendiente cuando se transmite esta respuesta.

Los caracteres \texttt{?} y \texttt{0x18} son tratados
por la rutina de recepción sin incorporarse a las líneas
G-code. El primero solicita un reporte de estado;
el segundo solicita el mensaje de bienvenida y
el restablecimiento de los offsets de trabajo.
Las acciones correspondientes se atienden mediante
banderas y la función \texttt{service\_realtime()},
desde el contexto del programa principal.

Durante los movimientos X--Y, esta función también
es invocada desde la espera de la rutina de movimiento.
De esta manera pueden atenderse solicitudes de estado
sin comenzar el procesamiento de otra línea G-code.

\subsection{Interpretación de las instrucciones}

Las instrucciones recibidas son cadenas de caracteres,
mientras que los módulos de control requieren valores
numéricos y acciones identificadas. Por ello se utiliza
un parser, es decir, una función que reconoce los
elementos de la instrucción y los transforma en una
representación que puede ser utilizada por el programa.

El análisis se implementa en \texttt{gcode\_parser.c}
mediante \texttt{gcode\_parse\_line()}. Su resultado
se almacena en una estructura \texttt{GcodeCommand},
que contiene el tipo de movimiento, los parámetros
numéricos y banderas que indican cuáles estaban
presentes en la línea.

Cada palabra se identifica mediante una letra y
un valor numérico. Por ejemplo, en
\texttt{G1X12.5Y8F1200S300}, la letra G identifica
el movimiento, X e Y indican coordenadas, F establece
el avance y S representa la consigna de accionamiento
del láser. Las banderas permiten distinguir un parámetro
omitido de uno cuyo valor es explícitamente cero.

El recorrido de la cadena admite letras mayúsculas
o minúsculas y palabras consecutivas sin espacios.
Se descartan los comentarios iniciados por punto
y coma y se omiten los comentarios entre paréntesis
en las líneas procesadas. Las líneas que comienzan,
después de los espacios iniciales, con un comentario
son descartadas.

La conversión numérica se realiza mediante
\texttt{parse\_gcode\_number()}, que reconoce un signo
opcional, dígitos y una parte decimal opcional.
Esta función evita interpretar la secuencia
\texttt{0X10} de una instrucción \texttt{G0X10}
como un número hexadecimal, lo que podría ocurrir
al utilizar una conversión numérica de propósito
general. Esta función fue implementada debido a una problemática
encontrada durante el proceso de prueba del firmware: el mensaje
\texttt{0x10} era interpretado como un número hexadecimal en vez de parte
de la orden GCode.

El parser no ejecuta movimientos ni modifica directamente
las salidas. La estructura obtenida es entregada a
\texttt{process\_gcode\_line()}, en \texttt{main.c},
donde se aplica el estado modal, se calculan los destinos
y se invocan las operaciones correspondientes.

\subsection{Comandos reconocidos y estado modal}

La implementación reconoce un subconjunto de instrucciones,
seleccionado para las funciones requeridas por el proyecto.
Este conjunto se resume en la Tabla~\ref{tab:gcode-admitido}.

\begin{table}[htbp]
    \centering
    \small
    \caption{Instrucciones y parámetros reconocidos
    por el firmware.}
    \label{tab:gcode-admitido}
    \begin{tabular}{ll}
        \hline
        Instrucción & Función \\
        \hline
        \texttt{G0} & Traslado con salida del láser anulada \\
        \texttt{G1} & Movimiento lineal con avance programado \\
        \texttt{G90}, \texttt{G91} & Coordenadas absolutas o relativas \\
        \texttt{G92}, \texttt{G92.1} & Establecer o borrar offsets X--Y \\
        \texttt{X}, \texttt{Y}, \texttt{Z} & Consignas de posición \\
        \texttt{F} & Avance expresado en mm/min \\
        \texttt{S} & Consigna de accionamiento del láser \\
        \texttt{M3}, \texttt{M4} & Habilitación constante o dinámica \\
        \texttt{M5} & Deshabilitación del láser \\
        \texttt{M17}, \texttt{M84} & Habilitación o liberación de motores \\
        \hline
    \end{tabular}
\end{table}

Los comandos de sistema se atienden directamente
en \texttt{main.c}. Entre ellos se encuentran
\texttt{\$H}, para solicitar la referenciación;
\texttt{\$X}, para limpiar la condición de alarma;
\texttt{\$J=}, para solicitar movimientos manuales;
y \texttt{\$\$}, para consultar la configuración
que el controlador informa al programa de computadora.

El modo de coordenadas y el avance programado se conservan
entre instrucciones. En los movimientos manuales
introducidos mediante \texttt{\$J=}, la selección
de coordenadas absolutas o relativas se aplica
a esa instrucción sin modificar el modo del programa.

Los offsets de trabajo permiten diferenciar las
coordenadas utilizadas por el trabajo de las posiciones
lógicas de máquina. En las consignas absolutas X--Y,
los offsets se suman a las coordenadas recibidas
para obtener el destino utilizado por el controlador.

Las instrucciones desconocidas detectadas por el parser
se identifican mediante la bandera \texttt{has\_unknown}.
En ese caso se rechaza la línea completa y se transmite
una respuesta de error. Este tratamiento no constituye
una implementación completa del lenguaje G-code ni
una validación exhaustiva de todas sus combinaciones.

\subsection{Anticipación de segmentos mediante lookahead}

Para reducir las detenciones entre segmentos consecutivos
se implementa una anticipación de un segmento,
denominada \textit{lookahead}. Una consigna X--Y válida
es almacenada en una estructura \texttt{PendingMove},
que conserva el destino, el avance, la dirección
normalizada y la información asociada al láser.

Cuando se recibe el siguiente destino X--Y válido,
se compara su dirección con la del segmento pendiente.
La comparación se realiza mediante el producto escalar
de ambos vectores normalizados:

\[
c = \hat{\mathbf{d}}_{\mathrm{pendiente}}
    \cdot \hat{\mathbf{d}}_{\mathrm{nuevo}}.
\]

Si se cumple $c > 0{,}999$, los segmentos se consideran
suficientemente alineados y con el mismo sentido.
Esta condición es utilizada para solicitar que
el segmento pendiente finalice a velocidad de crucero.
En caso contrario, se solicita su desaceleración final.

El segmento pendiente es ejecutado mediante
\texttt{flush\_pending\_move()}, y el nuevo segmento
queda almacenado para la comparación posterior.
Al calcular una nueva consigna relativa, se considera
como referencia el destino del pendiente, de modo
que las instrucciones mantienen su secuencia lógica.

Si no se recibe una continuación, el bucle principal
ejecuta el pendiente al comprobar que transcurrieron
más de 80~ms desde su almacenamiento. Esta condición
es evaluada por software y no corresponde a una
interrupción que se produzca exactamente al cumplirse
dicho tiempo.

También se ejecuta el pendiente antes de determinadas
operaciones que requieren conservar el orden,
como un movimiento Z, M5, M84, un cambio de offsets
o una actualización de S sin coordenadas X--Y.

La información del láser asociada al segmento es aplicada
al comenzar su ejecución. De esta manera, la recepción
de una consigna posterior no debe adelantar la aplicación
del valor S correspondiente a ese nuevo movimiento.

Este mecanismo conserva un solo segmento pendiente
y utiliza un criterio geométrico de continuidad.
No constituye un planificador de múltiples bloques
con optimización general de velocidades.

\subsection{Alcance de la compatibilidad con LaserGRBL}

El firmware implementa las funciones de comunicación
y los comandos necesarios para el uso previsto
con LaserGRBL. La consulta \texttt{\$\$} informa
los valores \texttt{\$30}, \texttt{\$31} y
\texttt{\$32}, asociados al intervalo de S y
a la identificación del modo láser.

La consigna máxima de S está configurada en 1000.
Este valor debe ser coherente con la configuración
utilizada al generar el trabajo. Las instrucciones M3
y M4 seleccionan accionamiento constante y dinámico,
respectivamente; el comportamiento del PWM se desarrolla
en la sección de control del láser.

La compatibilidad se limita a las funciones implementadas.
La transmisión de un mensaje de bienvenida con formato
GRBL no implica disponer de todas las instrucciones,
configuraciones o mecanismos de planificación del
firmware original. Por ello, el G-code utilizado
debe verificarse respecto del conjunto admitido.


% ============================================================================= %
\section{Generación y Coordinación del Movimiento}

\subsection{Cinemática del mecanismo CoreXY}

El desplazamiento del cabezal en el plano X--Y se
realiza mediante un mecanismo CoreXY. En esta
configuración, los motores A y B actúan sobre
una transmisión de correas que acopla sus movimientos
con las dos coordenadas cartesianas. Por lo tanto,
ninguno de los motores corresponde exclusivamente
al eje X o al eje Y.

La utilización de esta transformación en el firmware
responde a la configuración mecánica adoptada.
Las instrucciones G-code establecen destinos
cartesianos, mientras que los actuadores requieren
consignas de giro. En consecuencia, cada desplazamiento
solicitado debe ser convertido en los movimientos
correspondientes a los motores A y B.

Con la convención de signos utilizada en el proyecto,
la relación de movimiento se expresa como:

\begin{equation}
    \Delta a = \Delta x + \Delta y,
    \qquad
    \Delta b = \Delta x - \Delta y,
    \label{eq:corexy-inversa}
\end{equation}

donde $\Delta x$ y $\Delta y$ son los desplazamientos
cartesianos del cabezal, y $\Delta a$ y $\Delta b$
representan los desplazamientos equivalentes
de la transmisión asociados a los motores.
Todas estas magnitudes se expresan en unidades
de longitud.

La relación inversa permite obtener el desplazamiento
cartesiano a partir de la actuación de ambos motores:

\begin{equation}
    \Delta x = \frac{\Delta a + \Delta b}{2},
    \qquad
    \Delta y = \frac{\Delta a - \Delta b}{2}.
    \label{eq:corexy-directa}
\end{equation}

Estas relaciones muestran que un desplazamiento
sobre un único eje cartesiano requiere la actuación
de ambos motores. Para un movimiento exclusivamente
en X, ambos desplazamientos tienen igual magnitud
y signo. Para un movimiento exclusivamente en Y,
tienen igual magnitud y signos opuestos.

En cambio, cuando $\Delta x = \Delta y$, el
desplazamiento correspondiente al motor B es nulo.
Cuando $\Delta x = -\Delta y$, es nula la actuación
del motor A. Estos casos corresponden a diagonales
del plano cartesiano y permiten interpretar
el acoplamiento del mecanismo.

Los sentidos físicos de giro dependen de la disposición
de las correas, de la identificación de los motores
y de la conexión de sus bobinas. Por ello, la convención
utilizada por el firmware se ve reflejada en
el montaje de la máquina.

\subsection{Conversión de desplazamientos en pulsos}

La generación del movimiento requiere expresar
las consignas de los motores como cantidades
enteras de pulsos. Para ello se utiliza un factor
de conversión $K_s$, expresado en pulsos por milímetro,
que depende del ángulo de paso del motor,
la subdivisión de pasos seleccionada en el driver
y las características de la transmisión.

Para una polea dentada, el factor se calcula mediante:

\begin{equation}
    K_s =
    \frac{N_{\mathrm{rev}}\,m}{z\,p},
    \label{eq:pulsos-por-milimetro}
\end{equation}

donde $N_{\mathrm{rev}}$ es la cantidad de pasos
completos por revolución, $m$ es el factor de
subdivisión de pasos, $z$ es la cantidad de dientes
de la polea y $p$ es el paso de la correa.

En la configuración del firmware se consideran
motores de $1{,}8^\circ$ por paso, equivalentes
a 200 pasos completos por revolución, subdivisión
de $1/16$ de paso y poleas de 20 dientes con correas
de paso de $2\,\mathrm{mm}$. Con estos parámetros
se obtiene:

\begin{equation}
    K_s =
    \frac{200 \times 16}{20 \times 2}
    = 80\,\mathrm{pulsos/mm}.
\end{equation}

Este valor corresponde a los parámetros configurados
en el programa y debe ser contrastado con los
componentes y los ajustes utilizados en el montaje.

A partir del desplazamiento solicitado, las cantidades
de pulsos con signo se determinan como:

\begin{equation}
    N_A =
    \operatorname{round}
    \left[K_s(\Delta x + \Delta y)\right],
    \qquad
    N_B =
    \operatorname{round}
    \left[K_s(\Delta x - \Delta y)\right].
    \label{eq:pulsos-corexy}
\end{equation}

La operación de redondeo permite obtener valores
enteros compatibles con el accionamiento discreto
de los motores. El signo de cada resultado establece
el sentido de giro mediante la señal DIR, mientras
que su valor absoluto determina la cantidad
de pulsos que debe generarse en STEP.

En el módulo \texttt{stepper}, estas operaciones
se realizan a partir de la diferencia entre el
destino solicitado y la posición lógica conservada
por el firmware. La configuración permite, además,
invertir el sentido de X en la transformación
cinemática para adecuarlo al montaje.

\subsection{Interpolación y sincronización de los motores}

La transformación cinemática determina las cantidades
totales de pulsos necesarias para alcanzar el destino.
Sin embargo, estas cantidades no definen por sí solas
la trayectoria seguida durante el desplazamiento.
También debe establecerse cómo se distribuyen
los pulsos de ambos motores a lo largo del segmento.

Si todos los pasos de un motor fueran ejecutados
antes que los del otro, podría alcanzarse el mismo
destino ideal mediante una trayectoria diferente
de la recta solicitada. Por ello, los accionamientos
deben coordinarse durante la ejecución.

En el proyecto se utiliza una interpolación discreta
basada en el algoritmo de Bresenham. Se identifica
como accionamiento principal al motor que requiere
la mayor cantidad de pulsos:

\begin{equation}
    N_{\mathrm{mayor}} = \max(|N_A|,|N_B|),
    \qquad
    N_{\mathrm{menor}} = \min(|N_A|,|N_B|).
\end{equation}

En cada evento de temporización se genera un pulso
para el motor principal. La generación del pulso
del motor secundario se decide mediante un acumulador
de error, que permite distribuir sus pasos
proporcionalmente durante el recorrido.

El acumulador se inicializa con la mitad entera
de $N_{\mathrm{mayor}}$. En cada evento se resta
$N_{\mathrm{menor}}$; cuando el resultado es negativo,
se genera también un pulso para el motor secundario
y se suma $N_{\mathrm{mayor}}$ al acumulador.
De esta manera se conserva la proporción requerida
sin efectuar una división en cada decisión de paso.

Por ejemplo, para un desplazamiento de
$\Delta x = 10\,\mathrm{mm}$ y
$\Delta y = 5\,\mathrm{mm}$, con el factor
configurado de $80\,\mathrm{pulsos/mm}$, se obtiene:

\begin{equation}
    N_A = 1200,
    \qquad
    N_B = 400.
\end{equation}

El motor A constituye el accionamiento principal
y recibe un pulso en cada evento. Los 400 pulsos
del motor B se distribuyen durante el segmento,
aproximadamente uno cada tres eventos.
La relación entre ambos accionamientos permite
aproximar la trayectoria rectilínea mediante
movimientos discretos.

Cuando ambos motores requieren igual cantidad
de pulsos, se acciona cada uno en todos los eventos.
Cuando uno de ellos no requiere desplazamiento,
solo se generan pulsos para el otro.

La coordinación se realiza desde una misma rutina
de interrupción de TIM2. De esta manera,
las decisiones de paso de ambos motores comparten
una referencia temporal y un contador de avance
del segmento.

\subsection{Temporización y perfil de velocidad}

El parámetro F de la instrucción G-code representa
el avance a lo largo de la trayectoria, expresado
en milímetros por minuto. Para su utilización
en los cálculos de movimiento se convierte
a milímetros por segundo:

\begin{equation}
    v = \frac{F}{60}.
\end{equation}

La frecuencia de pasos del accionamiento principal
debe ajustarse según la orientación y la longitud
del segmento. Una misma frecuencia de motor no
produce necesariamente la misma velocidad cartesiana
en todos los movimientos CoreXY.

La longitud del desplazamiento solicitado se obtiene
mediante:

\begin{equation}
    L = \sqrt{(\Delta x)^2 + (\Delta y)^2}.
\end{equation}

Para una velocidad cartesiana nominal $v$, la
frecuencia de eventos del motor principal resulta:

\begin{equation}
    f_{\mathrm{mayor}} =
    v\,\frac{N_{\mathrm{mayor}}}{L}.
    \label{eq:frecuencia-corexy}
\end{equation}

En el firmware esta conversión se implementa
mediante un factor que relaciona los pulsos
del motor principal con la longitud cartesiana.
Así, el avance solicitado se interpreta respecto
del recorrido del cabezal y no de uno de los
accionamientos por separado.

TIM2 se configura con una base temporal nominal
de $1\,\mu\mathrm{s}$. Durante la ejecución,
su período se actualiza para establecer
el intervalo hasta el siguiente evento de paso.
La rutina de interrupción genera los pulsos,
actualiza el contador de avance y calcula
la temporización siguiente.

Para evitar cambios instantáneos de velocidad
se implementan fases de aceleración, velocidad
de crucero y desaceleración. Durante las rampas,
la velocidad de pasos se calcula a partir
de una relación de la forma:

\begin{equation}
    v_p^2 = 2a_p n,
\end{equation}

donde $v_p$ es la velocidad expresada en pulsos
por segundo, $a_p$ es la aceleración en pulsos
por segundo cuadrado y $n$ representa la cantidad
de eventos recorridos durante la aceleración
o restantes durante la desaceleración.

El período entre eventos se obtiene como
el inverso de la velocidad de pasos. También
se establece un período mínimo configurado
para limitar la frecuencia máxima de generación.

Las extensiones de las rampas se ajustan cuando
el segmento no dispone del recorrido previsto
para ambas fases. Asimismo, las condiciones
de entrada y salida a velocidad de crucero
pueden ser modificadas por el mecanismo
de \textit{lookahead}, descrito en la sección
anterior.

\subsection{Ejecución del barrido de grabado}

En el grabado por barrido, la imagen es representada
mediante una secuencia de recorridos lineales.
Durante cada línea se ejecutan segmentos consecutivos,
con consignas de accionamiento del láser asociadas
a la información que debe reproducirse sobre
el material.

Al finalizar una línea se realiza un desplazamiento
transversal para alcanzar la siguiente.
Según la configuración de generación de trayectorias,
el recorrido puede alternar su sentido entre líneas
o incluir un traslado para comenzar nuevamente
desde el mismo lado.

Las consignas cartesianas de cada segmento se
transforman mediante la cinemática CoreXY y sus
pulsos se distribuyen mediante la interpolación
descrita. Para las líneas paralelas a X,
ambos motores reciben cantidades iguales de pulsos
con el mismo signo; para el avance transversal
sobre Y, las cantidades son iguales y sus
signos opuestos.

El valor S puede variar entre segmentos de una
misma línea. Por lo tanto, la reproducción del
grabado requiere coordinar el desplazamiento
con la consigna del láser, que puede corresponder
a una salida nula o a diferentes ciclos de trabajo
del PWM.

En el modo M4, la implementación utiliza además
la relación entre la velocidad calculada por
el perfil y la velocidad de crucero para ajustar
el accionamiento del láser. Esta velocidad
es determinada por el programa; no corresponde
a una medición del movimiento físico.

\subsection{Movimiento independiente del eje Z}

El desplazamiento del eje Z se implementa
en el módulo \texttt{z\_axis} y se temporiza
mediante TIM4. Su consigna se convierte
en pulsos considerando los parámetros
del motor, el micropaso y el avance del husillo.

Los dos motores de Z se encuentran documentados
como conectados en paralelo a un mismo driver,
por lo que reciben una única secuencia
de accionamiento y no son controlados
individualmente por el firmware.

El eje Z no participa de la interpolación CoreXY
ni del almacenamiento de segmentos X--Y.
Su ejecución se realiza de manera independiente
para ajustar la posición vertical del cabezal.
Por ello, la implementación no debe interpretarse
como una interpolación simultánea de tres ejes.

\subsection{Resolución y exactitud del posicionamiento}

La conversión a cantidades enteras de pulsos
introduce una discretización del movimiento.
La subdivisión de pasos y la relación de transmisión
determinan la resolución nominal de las consignas,
pero no establecen por sí solas la exactitud
del desplazamiento físico.

La coordinación de los pulsos permite conservar
la relación de movimiento requerida por la
cinemática CoreXY y aproximar la trayectoria
programada. La exactitud resultante depende
también de la correspondencia entre los parámetros
configurados y el montaje real, de las holguras,
de la elasticidad y tensión de las correas
y de la ejecución de los pasos sin pérdidas.

El posicionamiento se realiza en lazo abierto.
Las posiciones conservadas por el firmware
son valores lógicos y no mediciones obtenidas
mediante sensores de posición. En la implementación
revisada, la posición lógica X--Y se actualiza
al destino antes de finalizar el desplazamiento.

Por lo tanto, el valor reportado por el programa
no constituye una verificación de que el cabezal
haya alcanzado físicamente esa posición.
La exactitud y la repetibilidad de la máquina
deben evaluarse mediante los ensayos correspondientes.


% ============================================================================= %
\section{Referenciación, Límites y Control del Láser}

% ============================================================================= %
\section{Ensayos y Resultados}

% ============================================================================= %
\section{Limitaciones y Conclusiones}

% ============================================================================= %

\begin{acknowledgments}
This research was supported by the Game Studies Research Initiative. We thank the anonymous reviewers for their valuable feedback and suggestions.
\end{acknowledgments}

%% References using BibTeX
%% Make sure to compile with: pdflatex -> bibtex -> pdflatex -> pdflatex
\bibliographystyle{plain}
\bibliography{references}


\end{document}
